\chapter{Diseño de evaluación exclusivamente NoSQL}

\section{Delimitación tecnológica}

Esta edición evalúa COBRIA únicamente dentro de sistemas NoSQL documentales. La unidad
de persistencia es el documento u objeto agregado; las consultas se expresan mediante
claves, filtros, índices del motor, colecciones derivadas y flujos de cambios. No se
transfieren cifras obtenidas con motores de otro paradigma. Esta decisión reduce el
volumen de evidencia cuantitativa disponible, pero alinea la inferencia con el objeto
real de investigación.

La evaluación distingue tres niveles. El nivel estructural verifica contratos en código
y configuración. El nivel ejecutable usa emuladores y entornos NoSQL aislados. El nivel
empírico mide latencia, costo, frescura y recuperación en condiciones reproducibles.
Una afirmación no asciende de nivel sin la evidencia correspondiente.

\section{Casos documentales anonimizados}

Se estudian tres aplicaciones con dominios distintos: eventos deportivos y multimedia,
operación hotelera y operación comercial. Sus nombres se omiten. Las tres emplean
documentos, colecciones, autenticación, funciones y almacenamiento administrado. Se
buscan contratos repetidos sin revelar organizaciones, personas ni credenciales.

\begin{table}[H]
\centering
\caption{Perfiles NoSQL utilizados en la evaluación.}
\begin{tabularx}{\textwidth}{p{2.4cm}X X}
\toprule
Perfil & Carga dominante & Riesgo arquitectónico observado\\
\midrule
N1 & eventos, resultados y multimedia & derivados obsoletos y reconstrucción costosa\\
N2 & reservas, disponibilidad y catálogos & concurrencia, tiempo y aislamiento\\
N3 & inventario, ventas y cierres & agregados operacionales y trazabilidad\\
\bottomrule
\end{tabularx}
\end{table}

\chapter{Hipótesis, variables y protocolo}

\section{Hipótesis activas}

\begin{enumerate}[leftmargin=*]
  \item[HN1.] Una proyección declarada por contrato puede reconstruirse desde documentos
  canónicos sin convertirse en autoridad paralela.
  \item[HN2.] El ámbito obligatorio impide recuperar documentos o fragmentos de otra
  organización antes de calcular similitud o construir contexto.
  \item[HN3.] Revisión, idempotencia y publicación por versión reducen duplicados,
  regresiones y estados parcialmente visibles.
  \item[HN4.] Un conjunto analítico versionado puede reproducirse desde fuentes,
  catálogos y transformaciones declaradas.
  \item[HN5.] Una proyección solo se justifica si mejora un objetivo de lectura y su
  mantenimiento permanece dentro de un presupuesto explícito.
\end{enumerate}

HN1--HN4 admiten verificación funcional. HN5 necesita mediciones específicas por motor,
región, índice, escala y mezcla de operaciones. Esta edición no publica un umbral
universal de lecturas por escritura.

\section{Variables y unidades}

Las variables independientes son motor documental, perfil, escala, número de ámbitos,
forma del documento, política de consistencia, estrategia de reconstrucción y nivel de
concurrencia. Las dependientes incluyen p50, p95 y p99; documentos y bytes leídos;
operaciones facturables; frescura; duplicados; divergencia; duración de reconstrucción;
violaciones de ámbito; y exactitud del conjunto analítico.

La corrida completa es la unidad inferencial. Cada manifiesto registra semilla, versión
de código, configuración, región, índices, calentamiento, orden de cargas y criterio de
descarte. Las enmiendas se publican antes de observar el resultado que podrían favorecer.

\chapter{Experimentos NoSQL}

\section{P1: lectura selectiva documental}

P1 compara la consulta canónica indexada con una proyección de cobertura sobre el mismo
conjunto lógico. Ambas configuraciones usan el mismo motor NoSQL y los mismos documentos
de entrada. La proyección gana únicamente si reduce el percentil objetivo o el costo
facturable sin alterar resultados, ámbito ni frescura más allá del límite acordado.

\section{P2: amplificación de escritura}

P2 actualiza documentos canónicos y registra operaciones adicionales, tamaño de eventos,
invocaciones, reintentos y tiempo hasta convergencia. El resultado es un vector de costo,
no una sola cifra. Una reducción de lectura no compensa automáticamente mayor complejidad.

\section{R1: reconstrucción total e incremental}

R1 elimina una versión derivada en un entorno controlado, reconstruye una candidata,
verifica equivalencia y publica mediante un cambio de versión. La variante incremental
procesa revisiones posteriores a un punto de control. Ambas rutas deben ser idempotentes,
conservar procedencia y producir el mismo resultado lógico.

\section{S1: aislamiento multiinquilino}

S1 genera identidades, documentos y consultas para múltiples ámbitos. Incluye ausencia
de ámbito, ámbito manipulado, claves de caché incompletas, búsqueda semántica y tareas
administrativas. La aceptación exige rechazo temprano y cero documentos cruzados. El
filtro posterior a la recuperación no satisface el contrato.

\section{A1: analítica e inteligencia artificial}

A1 reconstruye un conjunto fechado, publica variables versionadas y ejecuta recuperación
autorizada. Cada fragmento conserva fuente, ámbito, revisión, versión de transformación
y política. Se miden precisión, exhaustividad, abstención, deriva y fugas. La evaluación
de recuperación no se presenta como evaluación de un modelo generativo.

\chapter{Evidencia y resultados responsables}

\section{Evidencia disponible}

Los tres perfiles exhiben objetos canónicos, capas de acceso documental, servicios,
catálogos y representaciones derivadas. La recurrencia apoya la relevancia práctica del
problema, no la novedad de cada mecanismo. La contribución de COBRIA está en convertir
esa recurrencia en contratos comparables: autoridad, ámbito, revisión, tiempo, versión,
equivalencia y reconstrucción.

Las pruebas disponibles cubren reglas de acceso, matrices de roles, aislamiento, ciclos
de vida, idempotencia y reconstrucción. También existe evidencia longitudinal de menor
acceso directo desde presentación hacia la capa documental. Estos resultados prueban
rutas concretas, pero no productividad humana ni superioridad general.

\section{Resultados COBRIA Core 1.0}

La primera réplica congelada se ejecutó sobre dos motores documentales embebidos,
PouchDB y LokiJS, con escalas de 100, 1,000 y 5,000 documentos y treinta corridas por
celda. En total se conservaron 180 corridas. Todas las celdas superaron conjuntamente
P1, P2, R1, S1 y A1: equivalencia lógica completa, reconstrucción repetible,
amplificación de dos escrituras por cambio proyectado, rechazo de ámbito ausente y cero
documentos cruzados.

En PouchDB las medianas de lectura canónica fueron 1,186; 13,589 y 69,291 ms, frente a
0,440; 4,719 y 23,131 ms para la proyección. Las reconstrucciones medianas requirieron
8,376; 89,301 y 500,070 ms. En LokiJS las lecturas canónicas fueron 0,145; 1,263 y
7,026 ms, frente a 0,052; 0,409 y 2,220 ms; R1 requirió 0,663; 6,452 y 46,011 ms.

Estos resultados apoyan conformidad funcional y muestran una reducción de lectura en
este arnés embebido. No representan redes, regiones, concurrencia administrada, cuotas
ni facturación. HN1--HN4 quedan apoyadas para esta implementación de referencia; HN5
recibe apoyo local condicionado y permanece abierta para servicios administrados.

\chapter{Amenazas a la validez y síntesis}

Los casos comparten ecosistema y autoría. Los emuladores no reproducen cuotas,
particiones, latencia regional ni facturación administrada. Los datos anonimizados y
sintéticos pueden subrepresentar deriva y errores históricos. La inspección no sustituye
ejecución; las pruebas automáticas no sustituyen revisión humana; y una implementación
correcta no demuestra comprensión por todos los equipos.

La evidencia permite afirmar que COBRIA Core 1.0 es formalizable, ejecutable y conforme
en dos motores NoSQL documentales embebidos. No permite afirmar una mejora universal de
rendimiento ni un umbral económico transferible. El siguiente nivel de evidencia debe
provenir de servicios documentales administrados y una réplica externa.
